% ======================================================================
%  Section 4.3 and 4.3.1 --- Residual-based Tukey weighting.  Corrected.
%
%  The mathematics here is right. I checked rho, psi and w:
%      rho' = (c^2/6) * (-3)(1-u^2/c^2)^2 * (-2u/c^2) = u(1-u^2/c^2)^2
%      w = psi/u = (1-u^2/c^2)^2
%  and the claim that an overall factor cancels in psi(u)/u is correct.
%
%  YOUR %%% VERIFY IS ANSWERED: c = 4.6851 is confirmed against the
%  implementation. ViSP's vpRobust computes the rejection threshold as
%      double C = sigma * 4.6851;
%  so the constant in the thesis and the constant in the program agree, and
%  it is the same in Variants 2 and 3 as Section 4.6.2 requires. The marker
%  can be deleted.
%
%  FOUR CHANGES, one of which is substantive.
%
%  1. SUBSTANTIVE. The justification given for c = 4.6851 is 95 per cent
%     efficiency "when the residuals are Gaussian". The chapter elsewhere
%     says, twice, that they are not: 4.3.3 notes that gross misalignment
%     produces residuals that "do not follow the distribution the
%     M-estimator presumes", and 4.7.3 now measures the consequence, 13.6
%     per cent of valid measurements strongly down-weighted on a clean
%     scene at iteration 1. Quoting a Gaussian-model justification without
%     noting that the model does not hold here invites exactly the question
%     4.7.3 answers. One clause fixes it.
%  2. The forward reference for the rank failure pointed at the whole of
%     Section 4.2. It belongs to 4.2.2, where the mechanism and the
%     safeguard are now stated precisely.
%  3. "baseline" is used here for Variant 2 and in Remark 4.2.1 for
%     Variant 1. Two different things under one word.
%  4. u is called "the normalized residual" without being defined; and
%     -ize again.
% ======================================================================

\section{Residual-Based Tukey Weighting}
\label{sec:tukey}

%%% (3) Variant 1 is called "a genuine baseline" in Remark 4.2.1. Calling
%%% Variant 2 "the generic robust baseline" here gives the reader two
%%% baselines. "Reference" keeps them apart.
The first of the two weightings compared is the standard, purely
residual-based M-estimator, which serves in the ablation as the generic robust
reference against which the structure-aware proposal is measured.

\subsection{The Tukey biweight}
\label{sec:tukey-biweight}

Among the classical M-estimators, the Tukey biweight is adopted here because
it is \emph{redescending}: its influence function returns to zero for
sufficiently large residuals, so that measurements beyond a rejection
threshold are excluded from the estimate entirely rather than merely
down-weighted~\cite{key51}. This is the appropriate behaviour for occlusion,
where the corrupted measurements are not merely noisy but carry no information
about the pose at all. The associated loss function is
\begin{equation}
    \rho^{\mathrm{T}}(u)=
    \begin{cases}
        \dfrac{c^{2}}{6}\left[1-\left(1-\dfrac{u^{2}}{c^{2}}\right)^{3}\right],
            & \lvert u\rvert\le c,\\[10pt]
        \dfrac{c^{2}}{6}, & \lvert u\rvert>c,
    \end{cases}
    \label{eq:tukey-rho}
\end{equation}
where the saturation value $c^{2}/6$ follows from continuity at
$\lvert u\rvert=c$. Differentiating gives the influence function
$\psi^{\mathrm{T}}(u)=\rho^{\mathrm{T}\prime}(u)=u\bigl(1-u^{2}/c^{2}\bigr)^{2}$
on $\lvert u\rvert\le c$, the scale factors cancelling exactly. Conventions
differing by an overall non-zero multiplicative factor in $\rho^{\mathrm{T}}$
and $\psi^{\mathrm{T}}$ appear in the literature. They are immaterial for the
present IRLS formulation, since this factor cancels in the ratio $\psi(u)/u$
defining the weight~\eqref{eq:irls-weight}. Substituting into that expression
gives the Tukey weight function used throughout this chapter, understood as
the weight derived from the biweight loss~\eqref{eq:tukey-rho} rather than as
the loss or influence function themselves,
\begin{equation}
    w^{\mathrm{T}}(u)=
    \begin{cases}
        \left(1-\dfrac{u^{2}}{c^{2}}\right)^{2}, & \lvert u\rvert\le c,\\[10pt]
        0, & \lvert u\rvert>c,
    \end{cases}
    \label{eq:tukey}
\end{equation}
%%% (4) u is now defined where it is first used as a symbol in its own right.
where $u=e_{j}/\hat{s}$ is the residual normalised by the robust scale
estimate of Section~\ref{sec:mad}, and $c$ is the tuning constant governing
the rejection point.

%%% (1) The efficiency argument is kept, because c = 4.6851 is indeed the
%%% standard choice and was adopted for that reason. What is added is the
%%% acknowledgement that the model under which the figure is derived does
%%% not hold here, and a pointer to where the actual behaviour is measured.
%%% Without this the chapter asserts a Gaussian justification in 4.3.1 and
%%% denies the Gaussian assumption in 4.3.3 and 4.7.3.
The constant is set to $c=4.6851$, the standard choice for the biweight,
which yields approximately $95\%$ asymptotic efficiency relative to least
squares at the Gaussian model while retaining resistance to
contamination~\cite{key50}. It should be said that this justification is
asymptotic and rests on an assumption that does not hold in the present
setting: as Section~\ref{sec:structure-blind} argues and
Section~\ref{sec:weight-maps} measures, the photometric residuals of a
direct scheme under gross misalignment are neither independent nor Gaussian,
and a redescending estimator tuned at the Gaussian model rejects a substantial
fraction of uncorrupted measurements early in the servo. The value is
therefore adopted as the conventional one, held identical across Variants~2
and~3 so that the two differ in the argument supplied to the weight function
and not in the function itself, rather than because the model from which it
derives is satisfied here.

Note that the weight is exactly zero for $\lvert u\rvert>c$, not merely small.
%%% (2) was sec:robust-framework, the whole of 4.2. The mechanism and the
%%% safeguard both live in 4.2.2.
This hard rejection is precisely the mechanism identified in
Section~\ref{sec:weighted-law} as capable of reducing the rank of
$\mathbf{H}_{\mathbf{D}}$ and rendering the control law undefined, and it is
what makes the floor~\eqref{eq:hessian-floor} recorded there necessary rather
than precautionary.

% ======================================================================
%  LABELS THIS SECTION NOW NEEDS OR USES
%    sec:tukey-biweight   4.3.1, added above
%    sec:mad              4.3.2, referenced above -- add \label to it
%    sec:structure-blind  4.3.3, referenced above -- add \label to it
%    sec:weighted-law     4.2.2, added in the previous correction
%    sec:weight-maps      4.7.3
%    eq:hessian-floor     the floor added to (4.4) in the previous correction
%
%  EQUATION NUMBERING. Adding the floor as its own numbered equation in
%  4.2.2 shifts everything after it by one: what the current PDF prints as
%  (4.5) and (4.6) here become (4.6) and (4.7). LaTeX renumbers
%  automatically; the point is only that the numbers in the PDF you are
%  reading will no longer match after the rebuild.
%
%  NOT CHANGED
%  - The derivation is correct as it stands and I verified it term by term.
%  - The sentence about multiplicative conventions cancelling in psi(u)/u is
%    right and is worth keeping; it forestalls a reader who knows a
%    different normalisation of the biweight.
%  - cite key50 is used both for IRLS in 4.2.1 and for the efficiency figure
%    here. If key50 is a robust-statistics text that is fine; if it is a
%    paper specifically about IRLS, the efficiency claim needs its own
%    source. I cannot check the bibliography from the PDF.
% ======================================================================
